\section{Metodología Scrum + XP}
El equipo organiza el trabajo con Scrum y aplica prácticas de XP para construir y revisar el producto en incrementos. Se conserva la secuencia y el enfoque descritos en el informe del Sprint 2.

\subsection{Roles}
\begin{tabularx}{\textwidth}{|p{0.34\textwidth}|p{0.3\textwidth}|X|}\hline
\rowcolor{grisTarjeta}\textbf{Integrante} & \textbf{Rol} & \textbf{Responsabilidad}\\\hline
Jair León & Scrum Master / Backend & Facilita el trabajo del equipo y contribuye al backend.\\\hline
Paul Acurio & Product Owner Proxy / IA & Prioriza y aclara necesidades; apoya el análisis con IA.\\\hline
Alex Jonathan Chisag & Frontend / UI-UX & Implementa la interfaz y revisa la experiencia de uso.\\\hline
Rolando Sebastián Miranda & Backend / QA-DevOps & Apoya servicios backend, calidad e infraestructura.\\\hline
Ing. Hernán Naranjo, Mg. & Docente & Acompaña y revisa el trabajo académico.\\\hline
\end{tabularx}
\captionof{table}{Roles del equipo}

\subsection{Eventos y artefactos}
Se realizan planificación, reuniones diarias, revisión y retrospectiva. Los artefactos de trabajo son el Product Backlog, el Sprint Backlog, el incremento y la definición de hecho. El equipo revisa el avance y actualiza las tareas conforme a los acuerdos del sprint.

\subsection{Prácticas XP}
Se aplican pruebas automáticas, integración continua, programación en parejas o revisión de código, refactorización, entregas pequeñas, ramas de trabajo y solicitudes de cambios (PR). Estas prácticas ayudan a detectar errores durante el desarrollo y a integrar cambios revisables.

\subsection{Definición de Hecho}
Una historia se considera terminada cuando cumple sus criterios de aceptación, sus pruebas relevantes pasan, el cambio queda integrado y revisado, y la funcionalidad puede demostrarse con el incremento del sistema.
